\section{COMPARISON OF FUNCTIONS ON A FILTERED SET}

Let E be a set filtered by a filter with base F (Gen.~Top., I, p.~58); in this chapter we shall consider functions whose domain of definition is a subset of E belonging to the filter base F (the subset depending on the function) and which take their values either in the field of real numbers R, or, more generally, in a normed vector space over a valued field (Gen.~Top., IX, p.~170).

In applications, E will most often be a subspace of the real space \(R^{n}\) , or the extended line \(\overline{R}\) , and \(F\) will be the trace on E of the filter of neighbourhoods of a closure point of E, or else the filter of complements of relatively compact subsets of E (``neighbourhoods of the point at infinity'').

It will not in general suffice to know that such a function tends to a given limit along \(\mathfrak{F}\) to be able to treat all the problems of ``passage to the limit along \(\mathfrak{F}\)'' for expressions formed from this function.

For example, when the real variable \(x\) tends to \(+\infty\) the three functions \(x\) , \(x^2\) and \(\sqrt{x}\) all tend to \(+\infty\) , but, of the expressions

\[
(x + 1) ^ {2} - x ^ {2}, \qquad (x + 1) - x, \quad \sqrt {x + 1} - \sqrt {x},
\]

the first tends to \(+\infty\) , the second to 1, the third to 0.

It is thus important to know not just the limit value of a function along \(\mathfrak{F}\) (when this limit exists) but also the ``manner'' in which the function approaches its limit; in other words, one is led to classify the set of functions which tend to the same limit.
% semantic-fragments: legacy-input-seam
\relax{}
